%=====================================================================
%  CBEC-AI : IEEE Conference Paper (IEEEtran, two-column)
%  Compile on Overleaf: New Project -> Upload -> this .tex file
%  Menu -> Compiler: pdfLaTeX  (default). No extra files needed.
%=====================================================================
\documentclass[conference]{IEEEtran}
\IEEEoverridecommandlockouts

\usepackage{cite}
\usepackage{amsmath,amssymb,amsfonts}
\usepackage{algorithmic}
\usepackage{graphicx}
\usepackage{textcomp}
\usepackage{xcolor}
\usepackage{booktabs}
\usepackage{array}
\usepackage{multirow}
\usepackage{url}
\usepackage[T1]{fontenc}

% ---- convenience macros ----
\newcommand{\arrowto}{$\rightarrow$}

\begin{document}

\title{CBEC-AI: A Context-Aware AI-Driven Platform for\\
Cross-Border Export Intelligence, HS Classification,\\
and Regulatory Compliance for Indian SMEs}

\author{
\IEEEauthorblockN{<<AUTHOR~1 NAME>>, <<AUTHOR~2 NAME>>, <<AUTHOR~3 NAME>>}
\IEEEauthorblockA{\textit{<<Department of Computer Science / AI / IT>>} \\
\textit{<<Institution Name>>}\\
<<City>>, <<State>>, India \\
\{<<author1.email>>, <<author2.email>>, <<author3.email>>\}@<<domain>>}
}

\maketitle

%---------------------------------------------------------------------
\begin{abstract}
Small and Medium Enterprises (SMEs) constitute the backbone of India's
export economy yet face disproportionate barriers in cross-border trade:
complex Harmonized System (HS) code classification, fragmented
destination-country regulatory requirements, and the absence of
data-driven market selection tools. This paper presents \textbf{CBEC-AI},
an integrated, context-aware platform that unifies three AI-driven
capabilities: (i) a database-grounded, LLM-assisted HS code
classification engine that never fabricates codes; (ii) a hierarchical
regulatory intelligence retrieval system with \textbf{hybrid multimodal}
Retrieval-Augmented Generation (RAG)---combining structured
regulatory-database evidence with semantically retrieved ingested
documents, fused via a reranking stage---that answers compliance queries
strictly from official government evidence; and (iii) a machine-learning
export destination ranking model trained on eight years (2017--2024) of
bilateral trade data from CEPII BACI, CEPII Gravity, and World Bank WDI.
We evaluate the ranking model on a held-out temporal test set
(2023\arrowto 2024) across 55 HS6 product codes and India's top~10 export
destinations. After a leakage audit removed year-correlated features
(including an imputation-flag proxy and a synthetically year-formulated
GDP-per-capita feature) and models were retrained with time-aware
expanding-window cross-validation, a ranking-objective XGBoost model
(XGBRanker, \texttt{rank:pairwise}) achieves the best rank correlation
(Spearman~$\rho = 0.8956 \pm 0.0033$ across 5 seeds, Top-3 Overlap
$= 89.1\%$), significantly outperforming a naive historical-growth
baseline ($\rho = 0.8686$; Wilcoxon signed-rank $p = 3.8\times10^{-4}$)---in
contrast to an earlier, less rigorously feature-audited configuration in
which no model significantly beat the baseline. Random Forest remains the
strongest model for single-best-market selection (Top-1 Hit Rate
$= 81.8\%$), and a regression-objective XGBoost model remains best for
calibrated value prediction (RMSE $= 2.44$, $R^2 = 0.836$), indicating
that model choice should be matched to the downstream decision task rather
than a single ``best'' model declared across all metrics. The HS
classification engine achieves HIGH\_CONFIDENCE recommendations for
products in covered tariff chapters while honestly returning
NEEDS\_REVIEW for out-of-coverage products, and the served ranking model
returns DATA\_UNAVAILABLE outside its trained product/country scope rather
than extrapolating silently. To our knowledge, CBEC-AI is the first system
to combine grounded, hybrid multimodal LLM regulatory reasoning with a
leakage-audited, statistically validated ranking-objective trade model in
a single SME-facing platform, prioritizing correctness and auditability
over inflated performance claims.
\end{abstract}

\begin{IEEEkeywords}
Cross-border e-commerce, HS code classification, Retrieval-Augmented
Generation, export market ranking, trade facilitation, gradient boosting,
regulatory compliance, SME digitalization.
\end{IEEEkeywords}

%=====================================================================
\section{Introduction}
%=====================================================================

\subsection{Motivation}
India's merchandise exports crossed USD~450 billion in recent years, with
SMEs contributing approximately 45\% of this value. Despite this, SME
exporters operate at a significant information disadvantage compared to
large trading houses with dedicated compliance and market-research teams.
Three specific pain points recur:

\begin{enumerate}
\item \textbf{HS Code Misclassification.} The Indian Trade Classification
(ITC-HS) contains over 11{,}500 national tariff lines. Incorrect
classification results in customs delays, penalty duties, and cargo
detention. Manual classification requires specialized expertise
unavailable to most SMEs.
\item \textbf{Regulatory Fragmentation.} Each destination country imposes
distinct documentation, certification, labeling, and procedural
requirements that vary by product and HS code, scattered across dozens of
government portals (FDA, CBP, EU authorities, etc.).
\item \textbf{Suboptimal Market Selection.} SMEs frequently choose export
destinations based on anecdotal information rather than data-driven
analysis of demand, competition, tariffs, and market accessibility.
\end{enumerate}

\subsection{Contributions}
\begin{itemize}
\item \textbf{C1.} A \textbf{database-first HS classification
architecture} that combines deterministic multi-signal scoring with LLM
re-ranking, enforcing strict validation such that the AI can only rank
candidates retrieved from the official tariff database---eliminating
hallucination.
\item \textbf{C2.} A \textbf{grounded hybrid multimodal regulatory RAG
system} that fuses hierarchical HS matching (EXACT\_NATIONAL\_CODE
\arrowto{} HS6 \arrowto{} HS4 \arrowto{} HS2) with semantic document
retrieval and a multimodal reranking stage, transparently reports match
confidence, and refuses to fabricate requirements when evidence is absent.
\item \textbf{C3.} A \textbf{rigorously-evaluated export destination
ranking model} with honest statistical significance testing against naive
baselines, a 5-seed ablation study, and full reproducibility artifacts.
\item \textbf{C4.} An \textbf{end-to-end integrated platform} deployed as
a production web application connecting classification, compliance, and
market intelligence into a single SME workflow.
\item \textbf{C5.} A \textbf{reproducibility and integrity framework} that
separates verified data from AI-generated explanations and honestly
reports coverage gaps rather than inflating metrics.
\end{itemize}

\subsection{Design Philosophy}
The central architectural principle is \textit{``Database-first,
AI-second, honesty-always.''} The system never allows the language model
to be the source of regulatory or classification truth. Official
government data is authoritative; the AI ranks, retrieves, explains, and
summarizes but never invents. When data is unavailable, the system
explicitly says so.

%=====================================================================
\section{Related Work}
%=====================================================================

\subsection{Automated HS Classification}
Prior approaches include rule-based systems (WCO classification
databases), classical machine learning on customs declarations
\cite{ding2015,turhan2019}, and deep-learning text classifiers
\cite{kim2020}. A common failure mode is generating codes absent from the
official schedule. Recent LLM-based classifiers risk hallucinating
plausible-but-nonexistent codes. CBEC-AI addresses this by constraining
the LLM to rank only database-retrieved candidates and validating every
output against the tariff master.

\subsection{Trade Compliance Systems}
Commercial platforms (Descartes, E2open, Thomson Reuters ONESOURCE)
provide regulatory databases but are cost-prohibitive for SMEs, lack
transparent evidence lineage, and do not integrate AI explanation.
Academic work on regulatory NLP includes legal question answering
\cite{cui2023} and compliance checking \cite{zhang2023}, but few systems
apply grounded RAG to multi-jurisdictional trade compliance.

\subsection{Export Market Prediction and Gravity Models}
The gravity model of trade \cite{tinbergen1962,anderson2003} remains the
econometric standard for bilateral trade flows. Machine-learning
augmentations \cite{batarseh2021} show gains on trade volume prediction.
However, many studies omit rigorous significance testing against naive
baselines, risking overstated claims. Our work explicitly benchmarks
against a naive historical-growth baseline with paired non-parametric
tests.

\subsection{Retrieval-Augmented Generation}
RAG \cite{lewis2020} grounds LLM outputs in retrieved documents to
mitigate hallucination. We apply RAG specifically to trade regulatory
reasoning, with system-prompt constraints and explicit DATA\_UNAVAILABLE
fallbacks.

%=====================================================================
\section{System Architecture}
%=====================================================================

\subsection{Overview}
CBEC-AI comprises five layers: a React~19 single-page presentation layer;
a Spring~Boot~3.3 / Java~21 application layer hosting the core API, the HS
classification engine, and export analysis with compliance and RAG; an
AI~services layer (NVIDIA API: chat, embedding, multimodal reranking); a
dual-database persistence layer; and an offline data pipeline plus an
offline Python ML model. Fig.~\ref{fig:arch} summarizes the layering.

\begin{figure}[t]
\centering
\fbox{\parbox{0.92\columnwidth}{\footnotesize\ttfamily
PRESENTATION: React 19 + Vite (SPA)\\
\ \ Exporter Dashboard | Logistics | ML Ranking\\
\rule{\linewidth}{0.4pt}\\
APPLICATION: Spring Boot 3.3 / Java 21\\
\ \ [Core API] [HS Classify] [Export Analysis + Hybrid RAG]\\
\ \ [Market Opportunity] [Negotiation] [Compliance]\\
\rule{\linewidth}{0.4pt}\\
AI SERVICES: NVIDIA API\\
\ \ nemotron-3-ultra-550b (chat/explain)\\
\ \ llama-nemotron-embed-1b-v2 (embed)\\
\ \ llama-nemotron-rerank-vl-1b-v2 (rerank)\\
\rule{\linewidth}{0.4pt}\\
DATA: InternationalTrade DB | TradeData DB\\
\ \ auth/products/orders | 82{,}554 HS + regulations\\
\ \ document\_chunk (vector store)\\
\rule{\linewidth}{0.4pt}\\
OFFLINE ML: BACI+Gravity+WDI $\rightarrow$ features $\rightarrow$\\
\ \ XGBRanker/XGBRegressor/RF/LightGBM $\rightarrow$ ranking\\
PIPELINE (ETL): gov PDFs $\rightarrow$ parse $\rightarrow$ store
}}
\caption{CBEC-AI five-layer architecture. NVIDIA AI is a service layer
(not a data store); the reranker is a distinct stage between retrieval and
LLM generation. Offline ML includes the XGBRanker (v4 winner) alongside
regression models.}
\label{fig:arch}
\end{figure}

\subsection{Dual-Database Design}
A deliberate separation isolates transactional integrity from regulatory
data:
\begin{itemize}
\item \textbf{InternationalTrade DB:} users, products, orders, shipments
(Hibernate-managed, read-write).
\item \textbf{TradeData DB:} HS codes, regulations, documents,
certifications, HS-regulation mappings (read-only from backend, populated
exclusively by the ETL pipeline).
\end{itemize}
This ensures application bugs cannot corrupt official tariff data, and
regulatory updates deploy independently.

\subsection{Technology Stack}
\begin{table}[h]
\centering
\caption{Technology stack.}
\label{tab:stack}
\begin{tabular}{@{}ll@{}}
\toprule
\textbf{Layer} & \textbf{Technology} \\
\midrule
Frontend  & React 19, Vite 8, Tailwind CSS \\
Backend   & Spring Boot 3.3.0, Java 21, Spring Security + JWT \\
Databases & MySQL 8 (dual datasource via Hibernate) \\
AI        & NVIDIA API: nemotron-3-ultra-550b-a55b (chat), \\
          & llama-nemotron-embed-1b-v2 (embeddings), \\
          & llama-nemotron-rerank-vl-1b-v2 (multimodal rerank) \\
Pipeline  & Apache PDFBox, Jsoup, POI, OpenCSV \\
ML Model  & Python: XGBoost, LightGBM, scikit-learn, SHAP \\
\bottomrule
\end{tabular}
\end{table}

%=====================================================================
\section{HS Code Classification Engine}
%=====================================================================

\subsection{Pipeline}
The classification pipeline is: Product Input \arrowto{} Feature
Extraction \arrowto{} Candidate Retrieval \arrowto{} Deterministic Scoring
\arrowto{} LLM Re-ranking \arrowto{} Validation \arrowto{} Top-$N$.

\subsection{Feature Extraction}
Structured attributes are extracted and normalized: product name,
description, category, material, composition, function, manufacturing
process, physical form, and technical specifications. A domain terminology
map expands user vocabulary to HS nomenclature (e.g., ``cream'' \arrowto{}
\{preparation, beauty, skin care\}).

\subsection{Candidate Retrieval}
Three strategies execute without premature limiting: (1) direct
product-name LIKE search, (2) category-chapter keyword filtering, and (3)
MySQL FULLTEXT natural-language search. All unique candidates pass to
scoring; the top-20 are retained \emph{after} scoring---a critical design
decision preventing correct codes from being eliminated by alphabetical
truncation.

\subsection{Deterministic Scoring}
The composite score aggregates weighted signal matches:
\begin{align*}
S =\ & 0.25\,\text{descMatch} + 0.20\,\text{functionMatch} \\
     & + 0.15\,\text{materialMatch} + 0.10\,\text{categoryScore} \\
     & + 0.10\,\text{mfgMatch} + 0.10\,\text{techMatch} \\
     & + 0.05\,\text{physicalForm} + 0.05\,\text{officialDesc} \\
     & + \text{phraseBoost}(+40) + \text{hyphenBoost}(+35) \\
     & + \text{chapterAlign}(\pm 8/{-}18).
\end{align*}

\subsection{LLM Re-ranking}
The top-15 category-filtered candidates are sent to the LLM with a strict
prompt: \textit{``Rank ONLY the supplied candidates. Do NOT create,
modify, or invent an HS code.''} The final score is
$0.50\cdot\text{DB} + 0.50\cdot\text{AI}$. Every returned code is
validated to exist in \texttt{hs\_master} with \texttt{country='India'}
and \texttt{is\_current=true}.

\subsection{Confidence Levels}
\begin{table}[h]
\centering
\caption{Confidence level thresholds.}
\label{tab:conf}
\begin{tabular}{@{}ll@{}}
\toprule
\textbf{Score} & \textbf{Level} \\
\midrule
$\geq 90$ & HIGH\_CONFIDENCE \\
75--89    & GOOD\_MATCH \\
60--74    & POSSIBLE\_MATCH \\
40--59    & LOW\_CONFIDENCE \\
$<40$     & NEEDS\_REVIEW \\
\bottomrule
\end{tabular}
\end{table}

\subsection{Classification Results}
\begin{table}[h]
\centering
\caption{Representative classification outputs (current system:
\texttt{nemotron-3-ultra-550b-a55b} LLM re-ranking, AI\_RANKED mode).}
\label{tab:classres}
\begin{tabular}{@{}llrl@{}}
\toprule
\textbf{Product} & \textbf{HS Code} & \textbf{Score} & \textbf{Level} \\
\midrule
Cotton T-Shirt & 6109100000 & 96.2\%  & HIGH\_CONFIDENCE \\
Smartphone     & 8517130000 & 97.5\%  & HIGH\_CONFIDENCE \\
Basmati Rice   & 1006309059 & 77.8\%  & GOOD\_MATCH \\
Black Pepper   & 0904110020 & 75.9\%  & GOOD\_MATCH \\
Face Cream     & 3304991000 & 61.7\%  & POSSIBLE\_MATCH \\
\bottomrule
\end{tabular}
\end{table}
All five products are classified in AI\_RANKED mode with the current
\texttt{nemotron-3-ultra-550b-a55b} LLM. Compared to the earlier
\texttt{llama-3.1-8b-instruct} configuration, Smartphone's recommended
code shifted from 8517110000 to 8517130000 (a more specific national
tariff line) and Face Cream moved from 3304910000 to 3304991000 (a
narrower subheading); both remain within the correct HS4 heading.
For products in tariff chapters absent from the current dataset
(fertilizers, footwear, plastics), the system correctly returns
NEEDS\_REVIEW rather than hallucinating---a key integrity property.

%=====================================================================
\section{Regulatory Intelligence and RAG}
%=====================================================================

\subsection{Hierarchical Matching}
For a (destination country, HS code) query, the system matches at
progressively broader levels---EXACT\_NATIONAL\_CODE \arrowto{} HS6
\arrowto{} HS4\_HEADING \arrowto{} HS2\_CHAPTER---and always reports which
level was used, never presenting an HS4 match as product-specific.

\subsection{Data Model}
The TradeData database holds 103 regulations, 224 documents, 218
certifications, 221 labeling requirements, 105 restrictions, 110
procedures, and 389 HS-regulation mappings across 11 countries, sourced
from official authorities (FDA, USDA, EPA, FCC, CBP, EU customs, etc.).

\subsection{Grounded Hybrid Multimodal RAG}
The query flow is: Query \arrowto{} Retrieve (structured TradeData evidence
$+$ semantic vector search over ingested documents) \arrowto{} Rerank
\arrowto{} Build context \arrowto{} NVIDIA LLM (constrained) \arrowto{}
Grounded answer with citations. The system fuses two evidence streams: (a)
structured regulatory-database records from hierarchical HS matching, and
(b) semantically retrieved passages from ingested official documents
(e.g., 19~CFR~134 text, RN numbers, Textile Fiber Products Identification
Act provisions), embedded with \texttt{llama-nemotron-embed-1b-v2}. A
multimodal reranking stage \cite{nogueira2019}
(\texttt{llama-nemotron-rerank-vl-1b-v2})
orders the fused candidate set by relevance to the query before context
assembly. The system prompt enforces: \textit{``Use ONLY the supplied
regulatory context. Never invent regulations, certificates, authorities,
or dates. If evidence is insufficient, say so.''} When no evidence exists,
the response is \texttt{DATA\_UNAVAILABLE}. Degradation is graceful: if the
reranker is unavailable the fused candidates are used in retrieval order;
if the LLM is unreachable, structured regulatory data is still
returned---only the natural-language explanation is degraded.

\subsection{Export Analysis Example (India \arrowto{} USA, Black Pepper)}
\begin{table}[h]
\centering
\caption{Export analysis result, India \arrowto{} USA, black pepper.}
\label{tab:export}
\begin{tabular}{@{}lr@{}}
\toprule
\textbf{Metric} & \textbf{Value} \\
\midrule
Match Type            & HS4\_HEADING \\
Regulations           & 5 \\
Documents             & 14 \\
Certifications        & 13 \\
Labeling              & 15 \\
Restrictions          & 7 \\
Procedures            & 9 \\
Compliance Complexity & HIGH \\
\bottomrule
\end{tabular}
\end{table}

%=====================================================================
\section{Export Destination Ranking Model}
%=====================================================================

\subsection{Problem Formulation}
Given India's export profile for HS6 product $p$ in year $t$, rank
candidate destination countries by predicted export opportunity in year
$t{+}1$. The prediction target is $\log(\text{export value}_{t+1})$;
countries are ranked per product by predicted value.

\subsection{Dataset}
\begin{itemize}
\item \textbf{Products:} 55 HS6 codes across India's top 11 export
chapters (10, 27, 29, 30, 61, 62, 71, 72, 84, 85, 87).
\item \textbf{Countries:} India's top 10 export destinations by bilateral
trade volume (Bangladesh, China, Germany, Hong Kong, Netherlands, Saudi
Arabia, Singapore, UAE, UK, USA).
\item \textbf{Time:} 2017--2024 (8 years).
\item \textbf{Volume:} $55 \times 10 \times 8 = 4{,}400$ observations.
\item \textbf{Sources:} CEPII BACI (bilateral trade), CEPII Gravity
(distance, borders, FTAs), World Bank WDI (GDP, population).
\end{itemize}

\subsection{Feature Groups}
\begin{table}[h]
\centering
\caption{Feature groups.}
\label{tab:features}
\begin{tabular}{@{}p{1.7cm}p{6cm}@{}}
\toprule
\textbf{Group} & \textbf{Features} \\
\midrule
Trade      & export value, market share, import totals, HHI
             concentration, lags (1/2/3y), growth rates \\
Gravity    & distance, common border, common language, FTA/WTO \\
Macro      & destination GDP, GDP per capita, population \\
Engineered & demand per capita, export intensity, competition intensity,
             market-growth$\times$FTA \\
Product    & HS2/HS4/category encodings \\
Temporal   & year, trailing volatility, share change \\
\bottomrule
\end{tabular}
\end{table}

\subsection{Evaluation Protocol}
Held-out temporal split: 2023 features \arrowto{} 2024 realized exports.
Metrics: Spearman~$\rho$, Kendall~$\tau$, Top-3 Overlap, Top-1 Hit Rate,
NDCG@5/@10, RMSE, MAE, $R^2$. Global seed $=42$; ablations across 5 seeds
$\{42, 101, 2024, 777, 999\}$.

\subsection{Feature Audit and Leakage Correction}
\label{sub:audit}
Following the original ablation's discovery that an imputation-missingness
flag (\texttt{entry\_metrics\_imputed}) acted as an undisclosed proxy for
\texttt{year} ($r=0.873$), we conducted a systematic leakage audit \cite{kaufman2012} across
the full feature set prior to retraining. For every feature we computed
Pearson and Spearman correlation with \texttt{year}, mutual information
with \texttt{year}, and within-group (HS6~$\times$~year) variance, to
catch any feature that varies primarily across time rather than across the
product--country decision space it is intended to inform. The audit
removed three features:
\begin{itemize}
\item \textbf{\texttt{entry\_metrics\_imputed}}---confirmed as a binary
year-threshold indicator ($\mathbf{1}[\text{year}>2020]$), the same leak
identified in the prior ablation.
\item \textbf{\texttt{feature\_year\_num}}---the raw year value, which is
constant within each ranking group and falls outside the training range at
test time.
\item \textbf{\texttt{gdp\_per\_capita\_usd}}---removed per the prior
ablation finding that this feature degrades ranking performance
($p=0.032$) and was independently confirmed to carry a synthetic
year-trend.
\end{itemize}
We additionally identified that \texttt{gdp\_usd} and \texttt{population},
though sourced under the World Bank WDI label, are generated in the
data-build pipeline via a year-indexed formula rather than ingested
directly from WDI records. We retained these two features only for their
cross-sectional (country-level) signal and disclose this explicitly here;
we recommend regenerating them from authentic WDI extracts before further
publication or production use, and treat their presence as a limitation
(Section~\ref{sec:limitations}). The validated engineered
features---demand per capita, export intensity, competition intensity, and
market-growth~$\times$~FTA---were retained unchanged, consistent with
their statistically significant contribution ($p=0.0165$).

\subsection{Methodology Update: Time-Aware CV and Ranking-Objective Modeling}
All models were retrained on the leakage-corrected feature set using
expanding-window time-aware cross-validation (validation years 2019--2022
sequentially; the 2023\arrowto 2024 test split remained untouched
throughout development). In addition to the regression formulations, we
introduced an \texttt{XGBRanker} configuration (\texttt{rank:pairwise}
\cite{burges2010},
grouped by product~$\times$~year, with per-group integer relevance grades),
directly optimizing the pairwise ordering objective that the evaluation
metrics (Spearman~$\rho$, Top-$k$ Overlap, NDCG) measure, rather than a
proxy regression loss.

\subsection{Benchmark Results (v4)}
\label{sub:benchv4}
\begin{table*}[t]
\centering
\caption{Model benchmarking on leakage-corrected features (test:
2023\arrowto 2024; 5-seed mean$\pm$std for Spearman~$\rho$ where
applicable). Best value per column in bold.}
\label{tab:bench}
\begin{tabular}{@{}lrrrrrr@{}}
\toprule
\textbf{Model} & \textbf{Spearman $\rho$} & \textbf{Top-3} & \textbf{Top-1}
& \textbf{RMSE} & \textbf{$R^2$} & \textbf{Wilcoxon $p$} \\
\midrule
\textbf{XGBRanker (tuned)}   & $\mathbf{0.8956 \pm 0.0033}$ & \textbf{89.1\%} & 70.9\% & N/A$^{1}$ & N/A$^{1}$ & $\mathbf{3.8\times10^{-4}}$ \\
Random Forest                & $0.8775$ & 85.5\% & \textbf{81.8\%} & 2.47 & 0.83 & 0.042 \\
XGBoost Regressor (tuned)    & $0.8787$ & 85.5\% & 76.4\% & \textbf{2.44} & \textbf{0.836} & 0.114 \\
Naive Historical Growth      & $0.8686$ & 84.2\% & 74.5\% & 2.68 & 0.80 & --- \\
LightGBM                     & $0.8669$ & 83.0\% & 76.4\% & 2.52 & 0.825 & 0.762 \\
Gradient Boosting            & $0.8611$ & 81.8\% & 72.7\% & 2.57 & 0.819 & 0.854 \\
Ridge Regression             & $0.7135$ & 64.2\% & 41.8\% & 3.92 & 0.578 & $2.3\times10^{-7}$ \\
Heuristic Weighted           & $0.4555$ & 61.2\% & 65.5\% & N/A$^{1}$ & N/A$^{1}$ & $1.5\times10^{-10}$ \\
\bottomrule
\end{tabular}

\vspace{2pt}
{\footnotesize $^{1}$RMSE/$R^2$ are not applicable to XGBRanker (and the
heuristic), which produce uncalibrated relevance scores rather than point
value predictions. Spearman~$\rho$ for XGBRanker is the 5-seed mean$\pm$std;
other rows report the seed-42 point estimate (5-seed means for RF/XGB-Reg/
LightGBM/Naive differ by $<0.002$).}
\end{table*}

\textbf{Key finding (revised).} Unlike the original evaluation, in which the
best model's improvement over the naive baseline was not statistically
significant (Wilcoxon $p=0.124$), the ranking-objective XGBRanker achieves
a statistically significant improvement over the naive baseline
($p=3.8\times10^{-4}$) after feature-leakage correction and a
ranking-native training objective. Note that the naive baseline itself is
unaffected by the leakage audit---it uses only
\texttt{india\_export\_value\_usd} and \texttt{india\_export\_growth\_1y},
neither of which was removed---so its $\rho=0.8686$ is identical in both
v3 and v4. This indicates that part of the earlier
null result was attributable to (a)~residual leakage diluting the naive
baseline's apparent strength in relative terms, and (b)~a mismatch between
the regression training objective and the ranking evaluation metric,
rather than purely reflecting an intrinsic ceiling from trade inertia.
We emphasize, however, that \textbf{no single model dominates all metrics}:
best ranking order (Spearman~$\rho$, Top-3) is XGBRanker; best
single-destination selection (Top-1) is Random Forest; best calibrated
value estimate (RMSE, $R^2$) is the XGBoost Regressor. This
metric-dependent winner is itself evidence against a leakage explanation
for the XGBRanker's win: a leaked feature would be expected to inflate
performance uniformly across metrics, whereas here the model that wins on
ranking order is not the model that wins on point-value accuracy or top-1
selection.

Ridge Regression and the Heuristic baseline are decisively and
\textbf{significantly worse} than the naive baseline ($p < 0.01$ in both
cases), not better---their very low $p$-values (Table~\ref{tab:bench})
reflect the \emph{size} of the performance gap below the baseline, not an
advantage.

\textbf{Effect of leakage correction across model families.} The leakage
correction did not uniformly improve all models. LightGBM's $\rho$ moved
unfavorably ($0.8737 \rightarrow 0.8669$) and its Wilcoxon $p$-value
weakened ($0.346 \rightarrow 0.762$); Gradient Boosting's $p$-value also
weakened ($0.711 \rightarrow 0.854$) despite a small $\rho$ increase; and
Ridge's Top-1 dropped sharply ($49.1\% \rightarrow 41.8\%$). The
statistically significant win is specific to \emph{combining} leakage
correction \emph{with} a ranking-native training objective
(XGBRanker)---regression-style models did not uniformly benefit from
feature cleanup alone. This heterogeneous response further argues against
a simple ``removing bad features helps everyone'' interpretation and
supports the hypothesis that the objective-metric mismatch was an
independent contributing factor to the prior null result.

\subsection{Red-Flag Investigation}
Given the magnitude of the XGBRanker's improvement and its low $p$-value
relative to the original study, we treated the result as requiring
verification before acceptance, applying the same skepticism the original
ablation applied to \texttt{entry\_metrics\_imputed}. We confirmed:
\begin{enumerate}
\item \textbf{Stability across seeds}---$\rho = 0.8956 \pm 0.0033$ across 5
seeds $\{42,101,2024,777,999\}$, a variance profile similar to other
models.
\item \textbf{Significance against alternative models, not only the naive
baseline}---XGBRanker significantly outperforms the XGBoost Regressor
(paired $p=0.0020$) and Random Forest ($p=0.0023$) on the same feature set
and splits, ruling out a baseline-specific artifact.
\item \textbf{No leaked feature present}---the ranker's top features by
importance are conventional trade-momentum and gravity signals; none of
the three removed leakage features appear.
\item \textbf{Non-uniform performance across metrics}---the ranker's Top-1
Hit Rate (70.9\%) is \emph{below} both Random Forest and the naive
baseline, which would not occur if a leaked feature were driving the
result, since leakage typically inflates all correlated metrics
simultaneously.
\end{enumerate}
We report this investigation transparently rather than treating a favorable
$p$-value as sufficient justification on its own, consistent with the
integrity-first framing of this work.

\subsection{Ablation Study (Feature-Engineering, v3 Feature Set)}
Table~\ref{tab:ablation} documents the feature-engineering ablation on the
v3 feature set and is retained as historical reference. The v4 leakage
audit (Section~\ref{sub:audit}) is a distinct, additional correction layered
on top of these findings, not a replacement for them; re-running the full
5-seed ablation on the v4 feature set is noted as future work.
\begin{table*}[t]
\centering
\caption{5-seed ablation with paired $t$-tests vs.\ Full Model (df$=4$),
on the v3 feature set (historical reference).}
\label{tab:ablation}
\begin{tabular}{@{}lrrrrl@{}}
\toprule
\textbf{Configuration} & \textbf{Spearman $\rho$ (mean$\pm$std)} &
\textbf{Paired $\Delta\rho$} & \textbf{$t$} & \textbf{$p$} &
\textbf{Significant} \\
\midrule
Full Model (all features + year)  & $0.8679 \pm 0.0047$ & ---      & ---    & ---            & Reference \\
Model A: Trade only               & $0.8657 \pm 0.0026$ & $-0.0023$ & $-1.075$ & 0.343        & No \\
Model B: Trade + Macro            & $0.8629 \pm 0.0041$ & $-0.0050$ & $-3.975$ & \textbf{0.0165} & \textbf{Yes} \\
Model C: Full v1 (no engineered)  & $0.8631 \pm 0.0046$ & $-0.0048$ & $-2.043$ & 0.111        & No \\
Drop GDP per capita               & $0.8724 \pm 0.0036$ & $+0.0045$ & $+3.229$ & \textbf{0.032}  & \textbf{Yes} \\
Drop India Market Share           & $0.8672 \pm 0.0047$ & $-0.0007$ & $-0.239$ & 0.823        & No \\
Drop Distance (Gravity)           & $0.8656 \pm 0.0036$ & $-0.0023$ & $-2.317$ & 0.081        & No \\
\bottomrule
\end{tabular}
\end{table*}

\noindent\textbf{Findings.}
\begin{enumerate}
\item Omitting product-category and engineered intensity features
(Model~B) causes a \textbf{statistically significant} ranking decline
($p=0.0165$)---validating feature engineering.
\item Trade-only features (Model~A) achieve comparable correlation to the
full model ($p=0.343$), confirming trade momentum dominates.
\item Dropping GDP per capita \emph{improves} ranking significantly
($p=0.032$), indicating this macro feature adds noise for this task.
\end{enumerate}

\subsection{Data Quality Insight (v3 Feature Set)}
An investigation of the \texttt{entry\_metrics\_imputed} missingness flag
(for the discontinued World Bank Doing Business Index) revealed it ranked
\#5 in feature gain when the temporal trend was omitted. Adding an
explicit \texttt{year} feature dropped its importance to \#12
($-80.5\%$), confirming the flag was a \textbf{macro-period proxy}
($r=0.873$ with year) rather than a genuine ease-of-business
signal---a cautionary finding on leakage via missingness indicators.

\subsection{Two-Stage Hurdle Architecture (v3 Feature Set, Historical Reference)}
\label{sub:hurdle}
The following hurdle-model results were obtained on the v3 (pre-leakage-audit)
feature set and are retained here as a methodological reference. The hurdle
model was not re-evaluated on the v4 feature set; its absence from
Table~\ref{tab:bench} reflects this scope.

A zero-inflated hurdle model (Stage~1: zero-trade classifier
ROC-AUC$=0.9674$, PR-AUC$=0.9953$, accuracy$=94.18\%$; Stage~2:
positive-value regressor) achieves excellent zero-detection but does
\textbf{not} improve ranking over direct single-stage regression
($\rho=0.8599$ on v3 features), indicating the compound expectation
introduces variance without ranking benefit.

%=====================================================================
\section{Implementation and Deployment}
%=====================================================================

\subsection{Security}
BCrypt password hashing, JWT stateless tokens (24h expiry), account
lockout after 5 failed attempts, role-based access control (EXPORTER,
LOGISTICS), and environment-variable API-key management (NVIDIA keys never
exposed to frontend).

\subsection{Reliability Engineering for Frontier-Scale Inference}
The frontier-scale chat model is accessed through a resilience layer:
explicit connect/read timeouts, retry-with-backoff on transient
$5xx$/$404$/timeout errors, fail-fast on genuine $4xx$ errors (e.g.,
authentication), and fallback to structured-data-only responses if all
retries are exhausted. This ensures the compliance surface remains
available and never fabricates an answer when the LLM is degraded or
unreachable.

\subsection{API Surface}
Key endpoints: \texttt{/api/v1/hs/classify},
\texttt{/api/v1/export/analyze},
\texttt{/api/v1/market-opportunity/\{hs\}},
\texttt{/api/v1/regulations/\{country\}/\{hs\}},
\texttt{/api/v1/compliance/\{country\}/\{hs\}},
\texttt{/api/v1/ai/regulatory-chat}. The export-ranking endpoints
(\texttt{/api/v1/export/analyze} and
\texttt{/api/v1/market-opportunity}) are served by the v4 ranking model via
a precomputed prediction table, returning explicit \texttt{DATA\_UNAVAILABLE}
responses outside the model's 55-product~$\times$~10-country trained scope
rather than extrapolating silently.

\subsection{Reproducibility}
All ML experiments use fixed seeds. Raw per-seed values, benchmark
comparisons, ablation results, SHAP case studies, and data-quality reports
are persisted as JSON/CSV artifacts. Cross-table consistency is
programmatically asserted (Table~\ref{tab:bench} XGBoost $=$
Table~\ref{tab:ablation} Full Model seed 42 to $>4$ decimals).

%=====================================================================
\section{Discussion and Limitations}
\label{sec:limitations}
%=====================================================================
\begin{enumerate}
\item \textbf{HS Coverage:} The Indian HS dataset currently spans 21 of 99
chapters (extracted from the ITC-HS PDF). Out-of-coverage products
correctly receive NEEDS\_REVIEW. Completing coverage requires additional
official source ingestion.
\item \textbf{Regulatory Depth:} Most HS-regulation mappings are at HS4
level; the system honestly reports this rather than claiming exact
applicability.
\item \textbf{Ranking Signal:} After leakage correction and a
ranking-native objective, the XGBRanker significantly beats the naive
baseline ($p=3.8\times10^{-4}$), yet the absolute margin remains modest and
model choice is metric-dependent (Section~\ref{sub:benchv4}). Short-horizon trade
inertia still explains most of the predictable signal; we avoid
overclaiming and report per-metric winners rather than a single ``best''
model.
\item \textbf{Synthetic Macroeconomic Features:} During the v4 leakage
audit we identified that two features labeled as sourced from World Bank
WDI---\texttt{gdp\_usd} and \texttt{population}---are generated in the
data-build pipeline via a year-indexed formula rather than ingested
directly from authentic WDI extracts. GDP-per-capita, derived from these,
was removed entirely after being shown to degrade ranking in both the
original and revised ablations. \texttt{gdp\_usd} and \texttt{population}
were retained for their cross-sectional signal, but we flag this as a
data-provenance limitation: results involving these two features should be
treated as provisional pending regeneration from authentic WDI source
data. We disclose this rather than silently correcting or omitting it,
consistent with this paper's stated philosophy of honesty over polish.
\item \textbf{Country Scope:} The ranking model uses India's top-10 trade
partners; the deployed platform extends regulatory coverage to a broader
set using compliance scoring where ML training data is sparse.
\item \textbf{Ethical Positioning:} The system explicitly disclaims legal
authority (``verify with official customs authority'') and surfaces
confidence levels to prevent over-reliance.
\end{enumerate}

%=====================================================================
\section{Future Work}
%=====================================================================
\begin{itemize}
\item Complete 99-chapter Indian HS coverage via chapter-wise official
ingestion.
\item Product-specific (HS6/HS8) regulatory mappings using embedding
similarity.
\item Integration of official duty rates for landed-cost computation.
\item Multilingual product descriptions for rural exporters.
\item Longer prediction horizons and structural-break-aware models.
\item Federated learning from anonymized customs declarations.
\end{itemize}

%=====================================================================
\section{Conclusion}
%=====================================================================
CBEC-AI demonstrates a production-grade, integrity-first approach to
AI-assisted trade facilitation for Indian SMEs. Its defining
principle---database-first, AI-second, honesty-always---ensures no HS code
or regulatory requirement is ever fabricated, every recommendation is
traceable to official sources, and uncertainty is transparently
communicated. The export ranking evaluation contributes an honest,
rigorously-audited empirical finding: after removing year-correlated
leakage and matching the training objective to the ranking metric, a
ranking-objective model achieves a statistically significant---though
modest and metric-dependent---improvement over naive trade-momentum
extrapolation, while no single model dominates across ranking, top-1
selection, and calibrated value estimation. Equally important, the study
documents how a favorable result was interrogated for leakage before being
accepted. By prioritizing correctness, auditability, and honest reporting
over inflated metrics, CBEC-AI offers a trustworthy blueprint for deploying
AI in regulatory-sensitive domains.

%=====================================================================
\begin{thebibliography}{19}
%=====================================================================
\bibitem{lewis2020} P. Lewis \textit{et al.}, ``Retrieval-Augmented
Generation for Knowledge-Intensive NLP Tasks,'' in \textit{Proc. NeurIPS},
2020.
\bibitem{anderson2003} J. Anderson and E. van Wincoop, ``Gravity with
Gravitas: A Solution to the Border Puzzle,'' \textit{American Economic
Review}, 2003.
\bibitem{tinbergen1962} J. Tinbergen, \textit{Shaping the World Economy:
Suggestions for an International Economic Policy}. 1962.
\bibitem{ding2015} L. Ding \textit{et al.}, ``Automatic HS Code
Classification Using Machine Learning,'' in \textit{Proc. ICEC}, 2015.
\bibitem{kim2020} J. Kim \textit{et al.}, ``Deep Learning Approaches for
Customs Classification,'' \textit{IEEE Access}, 2020.
\bibitem{turhan2019} G. Turhan \textit{et al.}, ``Text Mining for
Harmonized System Classification,'' \textit{Expert Systems with
Applications}, 2019.
\bibitem{batarseh2021} F. Batarseh \textit{et al.}, ``Application of
Machine Learning in Forecasting International Trade Trends,''
\textit{Journal of Big Data}, 2021.
\bibitem{cui2023} J. Cui \textit{et al.}, ``Legal Question Answering with
Retrieval-Augmented LLMs,'' 2023.
\bibitem{zhang2023} Y. Zhang \textit{et al.}, ``Automated Compliance
Checking with Grounded Language Models,'' 2023.
\bibitem{baci2010} G. Gaulier and S. Zignago, ``BACI: International Trade
Database at the Product Level,'' \textit{CEPII Working Paper}, 2010.
\bibitem{cepii2022} M. Conte, P. Cotterlaz, and T. Mayer, ``The CEPII
Gravity Database,'' \textit{CEPII}, 2022.
\bibitem{wdi2024} World Bank, ``World Development Indicators,'' 2024.
\bibitem{wco2022} World Customs Organization, ``Harmonized System
Nomenclature 2022.''
\bibitem{xgboost2016} T. Chen and C. Guestrin, ``XGBoost: A Scalable Tree
Boosting System,'' in \textit{Proc. KDD}, 2016.
\bibitem{lightgbm2017} G. Ke \textit{et al.}, ``LightGBM: A Highly
Efficient Gradient Boosting Decision Tree,'' in \textit{Proc. NeurIPS},
2017.
\bibitem{shap2017} S. Lundberg and S. Lee, ``A Unified Approach to
Interpreting Model Predictions (SHAP),'' in \textit{Proc. NeurIPS}, 2017.
\bibitem{nogueira2019} R. Nogueira and K. Cho, ``Passage Re-ranking with
BERT,'' arXiv:1901.04085, 2019.
\bibitem{kaufman2012} S. Kaufman, S. Rosset, and C. Perlich, ``Leakage in
Data Mining: Formulation, Detection, and Avoidance,'' \textit{ACM TKDD},
vol.~6, no.~4, 2012.
\bibitem{burges2010} C. Burges, ``From RankNet to LambdaRank to
LambdaMART: An Overview,'' \textit{Microsoft Research Technical Report},
MSR-TR-2010-82, 2010.
\end{thebibliography}

%=====================================================================
\appendices
%=====================================================================
\section{Reproducibility Artifacts}
\begin{table}[h]
\centering
\caption{Reproducibility artifacts.}
\label{tab:artifacts}
\footnotesize
\begin{tabular}{@{}>{\raggedright\arraybackslash}p{2.9cm}>{\raggedright\arraybackslash}p{5.3cm}@{}}
\toprule
\textbf{Artifact} & \textbf{File} \\
\midrule
Master training dataset (4{,}400 rows) & \texttt{master\_cbec\_train\_dataset.csv} \\
Leakage audit (corr/MI vs year)        & \texttt{leakage\_audit\_v4.csv} \\
Selected features                      & \texttt{selected\_features\_v4.json} \\
Time-aware CV search results           & \texttt{cv\_results\_v4.csv} \\
Best hyperparameters                   & \texttt{best\_hyperparameters\_v4.json} \\
Benchmark (final)                      & \texttt{benchmark\_final\_v4.csv} \\
Seed-stability (5 seeds)               & \texttt{seed\_stability\_v4.csv} \\
Red-flag investigation                 & \texttt{redflag\_investigation\_v4.json} \\
Feature importance (ranker)            & \texttt{ranker\_feature\_importance\_v4.csv} \\
Final metrics summary                  & \texttt{final\_metrics\_summary\_v4.json} \\
Served prediction table                & \texttt{ml\_export\_rankings\_v4.csv} \\
Persisted models                       & \texttt{xgb\_ranker\_v4.json},
                                         \texttt{xgb\_regressor\_v4.json} \\
Training / investigation code          & \texttt{train\_models\_v4.py},
                                         \texttt{investigate\_v4.py} \\
Historical (v3) artifacts              & \texttt{benchmark\_comparison.csv},
                                         \texttt{ablation\_per\_seed\_raw.json} \\
\bottomrule
\end{tabular}
\end{table}

\section{Per-Seed Ablation Raw Values}
\begin{itemize}
\item \textbf{Full Model:} [0.8694, 0.8654, 0.8629, 0.8763, 0.8656],
mean$=0.8679$, std$=0.0047$.
\item \textbf{Model A (Trade Only):} [0.8626, 0.8681, 0.8626, 0.8686,
0.8664], mean$=0.8657$, std$=0.0026$.
\item \textbf{Model B (Trade+Macro):} [0.8631, 0.8647, 0.8592, 0.8694,
0.8581], mean$=0.8629$, std$=0.0041$.
\item \textbf{Model C (Full v1):} [0.8614, 0.8554, 0.8653, 0.8690,
0.8647], mean$=0.8631$, std$=0.0046$.
\end{itemize}

\end{document}
